\documentclass[10pt,a4paper]{amsart}
\usepackage[margin=19mm]{geometry}
\usepackage{amsmath,amssymb,mathtools}
\usepackage[scr=rsfs,cal=euler]{mathalfa}
\usepackage{xfrac}
\usepackage[unicode,hidelinks,bookmarks=false]{hyperref}
\newcommand{\ie}{i.e.\ }
\newcommand{\cf}{cf.\ }
\newcommand{\resp}{respectively\ }
\newcommand{\Subsection}{\subsection}
\providecommand{\nobreakdash}{\nobreak\hbox{-}}
\setlength{\emergencystretch}{2em}
\multlinegap0pt
\usepackage{amssymb}
\usepackage{mathtools}
\usepackage{bm}
\usepackage{xcolor}
\usepackage{float}
\usepackage[shortlabels]{enumitem}
\newtheorem{proposition}{Proposition}
\newtheorem{theorem}{Theorem}
\newcommand{\E}{\mathbb E}
\newcommand{\SCV}{\operatorname{SCV}}
\newcommand{\Var}{\operatorname{Var}}
\newcommand{\dd}{\,\mathrm d}
\newcommand{\e}{\mathrm e}
\title{\textbf{Least-Variable Harmonic Matrix-Exponential Distributions}}
\author{Maria Laura Battagliola, Oscar Peralta}
\date{Source: arXiv:2609.14612v1 (CC BY 4.0)}
\begin{document}
\maketitle

\noindent\emph{Source and licence.} \textbf{Least-Variable Harmonic Matrix-Exponential Distributions}. Version arXiv:2609.14612v1 (CC BY 4.0), distributed by arXiv under the Creative Commons Attribution 4.0 International licence, http://creativecommons.org/licenses/by/4.0/, as recorded in the arXiv licence metadata for this exact version and shown on the abstract page https://arxiv.org/abs/2609.14612v1. Full licence notice: \texttt{licenses/arxiv-2609.14612-CC-BY-4.0.txt}.\par\smallskip

\noindent\textit{Editorial scope.} Counted unit: the article's theorem thm:quadratic-construction (source0.tex lines 519--531). The article's complete original proof of it is delivered with it (source0.tex lines 533--598, one \texttt{proof} environment of 66 lines). No mathematics is written, repaired or re-derived: the statement and the proof are the article's own lines, in the article's own notation, and no retained line is substituted or re-flowed.  Retained uncounted as the source-local context the proof needs: lines 411--430; lines 446--448; lines 477--480.  Nothing was omitted from any retained span, so the package carries no omission record.  No retained line contains a citation, so the wrapper carries no bibliography and no bibliography entry.  No retained line contains a figure environment, an image-inclusion command or drawing code, so no image is delivered and the package declares no asset dependency.  Editorial formatting only, in addition: an AMS \texttt{amsart} preamble that restates the article's own declarations and macros used by the retained lines, namely the environments \texttt{proposition}, \texttt{theorem} on separate counters, together with the standard packages those lines need.  The retained environments print in this document as Proposition 1, Theorem 1 in order of appearance, whereas the article prints the same items with the numbering of its own full document.  The complete native source of the article as distributed in the arXiv e-print for this exact version is bundled unchanged as \texttt{sources/210337/source0.tex}.
\begin{proposition}\label{prop:period-decomposition}
Let $K$ be independent of $\Theta$ with geometric distribution
\[
  \Pr(K=k)=(1-q)q^k,
  \qquad k=0,1,2,\ldots.
\]
Then
\begin{equation}\label{eq:period-sum}
  Z\stackrel d=\Theta+2\pi K.
\end{equation}
Consequently,
\begin{align}
  \E[Z]
  &=\E[\Theta]+\frac{2\pi q}{1-q},
  \label{eq:period-mean}\\
  \Var(Z)
  &=\Var(\Theta)+\frac{(2\pi)^2q}{(1-q)^2}.
  \label{eq:period-variance}
\end{align}
\end{proposition}

\begin{equation}\label{eq:replica-negligible}
  L^2\frac{q_L}{(1-q_L)^2}\longrightarrow0.
\end{equation}

\begin{equation}\label{eq:circular-moment}
  \E\!\left[4\sin^2(U_m/2)\right]=\ell_m,
  \qquad m^2\ell_m\longrightarrow\pi^2.
\end{equation}

\begin{theorem}[Quadratic harmonic construction]\label{thm:quadratic-construction}
For every fixed $\vartheta\in(\pi,2\pi)$, the preceding construction satisfies
\begin{equation}\label{eq:construction-moments}
  \E[Z_{L,\vartheta}]\longrightarrow\vartheta,
  \qquad
  \limsup_{L\to\infty}L^2\Var(Z_{L,\vartheta})\leq\pi^2.
\end{equation}
Consequently,
\begin{equation}\label{eq:construction-bound}
  \limsup_{L\to\infty}(2L+1)^2\SCV(Z_{L,\vartheta})
  \leq\frac{4\pi^2}{\vartheta^2}.
\end{equation}
\end{theorem}

\begin{proof}
Write $m=m_L$, $r=r_L$ and $\beta=\beta_L$, and let $\Theta_{L,\vartheta}$ denote the first-period component in Proposition~\ref{prop:period-decomposition}. We bound its second moment about $\vartheta$, which also bounds its variance.

Use $u\in[-\pi,\pi)$ as the representative of $z-\vartheta$ modulo $2\pi$, and put $w_L(u)=h_m(u)\cos^{2r}(u/2)$. For $u<\delta$ the position in the first period is $z=\vartheta+u$, and for $u\geq\delta$ it is $z=\vartheta+u-2\pi$. After cancelling the common factors, the normalizing constant is therefore
\[
  D_L=\int_{-\pi}^{\delta}\e^{-\beta u}w_L(u)\dd u
        +L^3\int_{\delta}^{\pi}\e^{-\beta u}w_L(u)\dd u.
\]
The factor $L^3=\e^{2\pi\beta}$ accounts for the part of the circular profile that wraps into the beginning of the first period. The normalizer $D_L$ is kept explicit because the following estimates first control unnormalized second-moment contributions. The lower bound $D_L\geq1-o(1)$ established below ensures that normalization does not increase the leading asymptotic upper constant.

We first bound $D_L$ from below, using the elementary inequality $1-(1-x)^r\leq rx$ for $0\leq x\leq1$, which gives
\[
  \E\!\left[1-\cos^{2r}(U_m/2)\right]\leq\frac{r\ell_m}{4}.
\]
Also, by \eqref{eq:circular-moment},
\[
  \Pr(|U_m|\geq\delta)
  \leq\frac{\ell_m}{4\sin^2(\delta/2)}.
\]
Since $w_L$ is even and $[-\delta,\delta]\subset[-\pi,\delta]$, pairing the contributions at $u$ and $-u$ replaces the exponential weights by $\cosh(\beta u)\geq1$, and we obtain
\begin{align}
  D_L
  &\geq\int_{-\delta}^{\delta}\cosh(\beta u)w_L(u)\dd u
   \geq\int_{-\delta}^{\delta}w_L(u)\dd u \notag\\
  &\geq1-\frac{r\ell_m}{4}
          -\frac{\ell_m}{4\sin^2(\delta/2)}
   =1-o(1).
  \label{eq:normalizer-bound}
\end{align}

Next, consider the second moment on the part that does not wrap. Here, the potentially increasing factor $\e^{\beta|u|}$ from the exponential tilt competes with the Gaussian-type decay supplied by the cosine power. The exponent $\beta|u|-(r-1)u^2/4$ below records exactly this competition. For $|u|\leq\pi$,
\begin{equation}\label{eq:elementary-cosine-bounds}
  u^2\cos^2(u/2)\leq4\sin^2(u/2),
  \qquad
  \cos(u/2)\leq\e^{-u^2/8}.
\end{equation}
The first inequality follows from $x\leq\tan x$ for $0\leq x<\pi/2$. Integrating that inequality gives $-\log\cos x\geq x^2/2$, and hence the second, with the endpoint values following by continuity. As $r>1$, completing the square yields
\begin{align*}
  u^2\e^{-\beta u}\cos^{2r}(u/2)
  &\leq4\sin^2(u/2)
      \exp\!\left(\beta|u|-\frac{r-1}{4}u^2\right)\\
  &\leq4\sin^2(u/2)\exp\!\left(\frac{\beta^2}{r-1}\right).
\end{align*}
Integration against $h_m$ thus bounds the unnormalized second-moment contribution from $u<\delta$ by $\ell_m\exp(\beta^2/(r-1))$.

On the wrapped part, $u\geq\delta>0$ and the displacement from $\vartheta$ is $u-2\pi$. Every wrapped point is therefore at least circular distance $\delta$ from the peak, and $\cos^{2r}(u/2)\leq\exp(-ru^2/4)\leq\exp(-r\delta^2/4)$. This exponential suppression is what defeats the $L^3$ amplification caused by wrapping. Therefore
\[
  L^3\int_{\delta}^{\pi}(u-2\pi)^2\e^{-\beta u}w_L(u)\dd u
  \leq4\pi^2L^3\e^{-r\delta^2/4}.
\]
Here, we used $(u-2\pi)^2\leq4\pi^2$, $\e^{-\beta u}\leq1$, \eqref{eq:elementary-cosine-bounds}, and $\int_{\delta}^{\pi}h_m\leq1$. Since $r$ is of order $\sqrt L$ and $\delta>0$ is fixed, $L^3\e^{-r\delta^2/4}=o(L^{-2})$. The exponential decay in $\sqrt L$ dominates every polynomial factor. Combining the two contributions gives
\[
  \E\!\left[(\Theta_{L,\vartheta}-\vartheta)^2\right]
  \leq\frac{\ell_m\exp(\beta^2/(r-1))
                  +4\pi^2L^3\e^{-r\delta^2/4}}{D_L}.
\]

Now $m/L\to1$, $L^2\ell_m\to\pi^2$, and $\beta^2/(r-1)\to0$. Because $\delta$ is fixed and positive, the wrapped bound is $o(L^{-2})$. Together with \eqref{eq:normalizer-bound}, this proves
\[
  \limsup_{L\to\infty}
  L^2\E\!\left[(\Theta_{L,\vartheta}-\vartheta)^2\right]\leq\pi^2.
\]
In particular, $\E[\Theta_{L,\vartheta}]\to\vartheta$ by Cauchy--Schwarz, and the same upper bound holds for $L^2\Var(\Theta_{L,\vartheta})$.

Finally, Proposition~\ref{prop:period-decomposition} adds an independent geometric number of periods with parameter $q_L=L^{-3}$. Its mean contribution vanishes and its variance contribution is $o(L^{-2})$ by \eqref{eq:replica-negligible}. This proves \eqref{eq:construction-moments}. Dividing the variance by the squared mean and using $(2L+1)^2/L^2\to4$ gives \eqref{eq:construction-bound}.
\end{proof}

\end{document}
